
\section{Mitigación aplicada por Cloudflare}

Para recuperar el servicio, los ingenieros de Cloudflare
utilizaron procedimientos de acceso de respaldo, ya que
el fallo también dificultaba la conexión con los centros
afectados. Una vez identificada la causa, comenzaron a
revertir la configuración incorrecta de los routers.
La recuperación fue gradual y, a las 07:42 UTC, los
19 centros de datos afectados estaban nuevamente
operativos \cite{cloudflare2022}.

Después del incidente, Cloudflare anunció mejoras en
sus procedimientos. Entre ellas, realizar pruebas en
centros con arquitectura MCP antes de extender los
cambios, revisar las reglas de enrutamiento para evitar
errores en su orden y automatizar la reversión de
configuraciones cuando se detecten problemas.
Estas medidas buscan reducir el impacto de futuros
fallos y acelerar la recuperación del servicio
\cite{cloudflare2022}.
